% xelatex -jobname=statement main.tex (run twice for page references)
\documentclass[12pt,a4paper,oneside]{ctexart}
\usepackage{xeCJKfntef,xcolor}
\usepackage{tabularx}
\usepackage{array}
\usepackage{bookmark}
\usepackage{parallel}
\usepackage{clrscode3e}
\usepackage{lastpage}
\usepackage{amsfonts}
\usepackage{amssymb}
\usepackage{amsmath}
\usepackage{bm}
% No external syntax highlighter is needed; samples use verbatim.
\usepackage{hyperref}
\usepackage{graphicx}
\usepackage{enumitem}
\usepackage{multirow}
\usepackage{float}
\usepackage{booktabs}
\hypersetup{
colorlinks=true,
linkcolor=blue,
citecolor=blue
}
\usepackage[english]{babel}
\usepackage{setspace}
\usepackage{geometry}
\geometry{left=2cm,right=2cm,top=2cm,bottom=2cm}
\usepackage{indentfirst}

\setlength{\parindent}{2em}
\emergencystretch=2em
\linespread{1.15}
\setcounter{secnumdepth}{0}

\newcommand{\bfemph}[1]{\CJKunderdot{\textbf{#1}}}
\renewcommand{\emph}[1]{\bfemph{#1}}
\setlength{\headheight}{18pt}
\setlist[itemize]{labelsep = 1.1em, itemsep = 0pt, parsep = 1pt, itemindent = 1.5em, topsep = 0pt}
\setlist[enumerate]{labelsep = 1.1em, itemsep = 0pt, parsep = 1pt, itemindent = 1.5em, topsep = 0pt}
\newcommand{\tightlist}{}

\usepackage{fancyhdr}
\pagestyle{fancy}
\renewcommand{\headrulewidth}{0.4pt}
\renewcommand{\footrulewidth}{0.4pt}
\setlength{\headheight}{18pt}
\lhead{NOIP 模拟赛}
\chead{}
\rhead{\rightmark}
\lfoot{}
\cfoot{\small 第 \thepage\ 页  \hspace{2em} 共 \pageref{LastPage} 页}
\rfoot{}

\newcommand{\problemtitle}[2]{%
  \section*{#1}%
  \markright{#1}%
  \phantomsection\label{#2}%
  \pdfbookmark[1]{#1}{#2}%
}
\newcommand{\problemheading}[2]{%
  \subsection*{#1}%
  \markright{#1}%
  \phantomsection\label{#2}%
  \pdfbookmark[2]{#1}{#2}%
}
\newenvironment{centeredtable}{\par\vspace{.35em}\setlength{\tabcolsep}{4pt}\begin{center}}{\end{center}\vspace{.35em}}

\begin{document}
  \title{\heiti\zihao{1} NOIP 模拟赛}
  \author{ASDFZ}
  \date{2026年10月24日}
  \maketitle
  \begin{center}
    \footnotesize
    \begin{tabularx}{\textwidth}{|X|X|X|X|X|}
    \hline
题目名称 & 叶环独立集 & 窗口校准 & 上下左右 & 存储空间\\
\hline
题目类型 & 传统型 & 传统型 & 传统型 & 传统型\\
\hline
目录 & \texttt{ring} & \texttt{calibration} & \texttt{lrud} & \texttt{storage}\\
\hline
可执行文件名 & \texttt{ring} & \texttt{calibration} & \texttt{lrud} & \texttt{storage}\\
\hline
输入文件名 & \texttt{ring.in} & \texttt{calibration.in} & \texttt{lrud.in} & \texttt{storage.in}\\
\hline
输出文件名 & \texttt{ring.out} & \texttt{calibration.out} & \texttt{lrud.out} & \texttt{storage.out}\\
\hline
测试点时限 & 1 秒 & 1 秒 & 1 秒 & 10 秒\\
\hline
内存限制 & 128 MiB & 512 MiB & 512 MiB & 256 MiB\\
\hline
    \end{tabularx}
  \end{center}
  \hspace{2em}提交源程序文件名
  \begin{center}
    \footnotesize
    \begin{tabularx}{\textwidth}{|p{90pt}|X<{\centering}|X<{\centering}|X<{\centering}|X<{\centering}|}
    \hline
    对 C++ 语言 & \texttt{ring.cpp} & \texttt{calibration.cpp} & \texttt{lrud.cpp} & \texttt{storage.cpp}\\
    \hline
    \end{tabularx}
  \end{center}
  \hspace{2em}编译选项
  \begin{center}
    \begin{tabularx}{\textwidth}{|p{100pt}|X<{\centering}|}
    \hline
    对 C++ 语言 & \texttt{-O2 -std=c++17} \\
    \hline
    \end{tabularx}
  \end{center}
  \vspace{.8em}
  \hspace{1em}{\heiti 注意事项与提醒：}
  \begin{enumerate}[topsep = 4pt]
    \setlength{\labelsep}{1.1em}
    \setlength{\itemsep}{2pt}
    \setlength{\parsep}{2pt}
    \setlength{\parskip}{2pt}
    \setlength{\itemindent}{.5em}
    \item 文件名（程序名和输入输出文件名）必须使用英文小写。
    \item C/C++ 中函数 main() 的返回类型必须是 int，程序正常结束时的返回值必须为 0。
    \item 第三题（lrud）采用特殊评测，其余题目采用全文比较（过滤行末空格及文末回车）。
    \item 程序可使用的栈内存空间限制与题目的内存限制一致。
    \item 只提供 Linux 格式的附加样例文件。
  \end{enumerate}
  \thispagestyle{empty}
  \newpage
% BEGIN GENERATED PROBLEMS
\problemtitle{叶环独立集(ring)}{probA_title}

\noindent 时间限制：1 秒\qquad 空间限制：128 MiB\par

\problemheading{题目描述}{probA_description}

给定一棵有根树，共有 \(n\) 个顶点，编号为 \(1,2,\ldots,n\)，根为 \(1\)。对于每个 \(i>1\)，顶点 \(i\) 的父亲是 \(p_i\)，并且 \(p_i<i\)。没有孩子的顶点称为\textbf{叶子}。

从根开始进行深度优先遍历。遍历一个顶点时，按照孩子编号从小到大的顺序访问它的孩子。设叶子按照此次 DFS 中的访问顺序依次为 \(l_1,l_2,\ldots,l_k\)。

现在构造一个无向图 \(G\)：

\begin{itemize}
\tightlist
\item
  对于树中的每条父子边，在 \(G\) 中连接对应的两个顶点。
\item
  对于叶子序列中相邻的每一对 \(l_i,l_{i+1}\)（\(1\le i<k\)），在 \(G\) 中连接这两个叶子。
\item
  连接 \(l_k\) 和 \(l_1\)。
\end{itemize}

如果新增边重复，则只保留一条；不添加自环。因此，只有一个叶子时不添加叶子边，只有两个叶子时它们之间只有一条边。

\textbf{独立集}是一个顶点集合，其中任意两个不同顶点都不相邻。

请计算图 \(G\) 的独立集总数（包括空集），并对 \(1\,000\,000\,007\) 取模。

\problemheading{输入格式}{probA_input}

第一行包含一个整数 \(n\)。

第二行包含 \(n-1\) 个整数 \(p_2,p_3,\ldots,p_n\)，其中 \(p_i\) 表示顶点 \(i\) 的父亲。

\problemheading{输出格式}{probA_output}

输出一个整数，表示独立集总数对 \(1\,000\,000\,007\) 取模后的结果。

\par\medskip\noindent\begin{minipage}{\linewidth}
\problemheading{样例 1 输入}{probA_sample_1_input}

\begin{verbatim}
7
1 1 2 2 3 3
\end{verbatim}
\end{minipage}\par\medskip

\par\medskip\noindent\begin{minipage}{\linewidth}
\problemheading{样例 1 输出}{probA_sample_1_output}

\begin{verbatim}
21
\end{verbatim}
\end{minipage}\par\medskip

\problemheading{样例 1 解释}{probA_sample_1_explanation}

树的叶子依次为 \(4,5,6,7\)，它们在图中构成一个四边形。结合树边限制，图的独立集总数为 \(21\)。

\problemheading{数据规模与约定}{probA_constraints}

所有数据均满足 \(4\le n\le 10^6\)，且 \(1\le p_i<i\)。树中允许出现只有一个孩子的顶点。

本题采用分组计分。一个子任务内全部测试点通过后，才能获得该组分数；各组分数独立相加。

\begin{centeredtable}\begin{tabularx}{0.85\textwidth}{|c|c|X|}\hline
子任务 & 分值 & 额外限制\\\hline
1 & 15 & \(n\le 20\)\\\hline
2 & 25 & \(n\le 10^5\)，叶子数不超过 \(8\)\\\hline
3 & 15 & 这棵树是以 \(1\) 为中心的星形树\\\hline
4 & 45 & 无额外限制\\\hline\end{tabularx}\end{centeredtable}
\newpage
\problemtitle{窗口校准(calibration)}{probB_title}

\noindent 时间限制：4 秒\qquad 空间限制：512 MiB\par

\problemheading{题目描述}{probB_description}

给定长度为 \(n\) 的序列 \(a_1,a_2,\ldots,a_n\)。序列中的元素均为模 \(P=998244353\) 意义下的整数，以下所有运算也均在模 \(P\) 意义下进行。

你需要依次处理 \(m\) 次操作。

\textbf{操作一：\texttt{1\ l\ r\ k}，校准窗口。}

保证 \(1\le k\le r-l\)。对区间 \([l,r]\)，定义长度为 \(k\) 的窗口和

\[W_i=\sum_{j=i}^{i+k-1}a_j\qquad(l\le i\le r-k+1).\]

保留 \(a_l,a_{l+1},\ldots,a_{l+k}\) 的当前值，修改 \(a_{l+k+1},\ldots,a_r\)，使这些窗口和按起点递增的顺序构成等差数列。也就是说，存在一个模 \(P\) 意义下的整数 \(c\)，满足

\[W_{i+1}-W_i=c\qquad(l\le i\le r-k).\]

这里的 \(W_i\) 均由修改后的序列计算，\(c\) 可以在不同操作中取不同值。区间 \([l,r]\) 外的元素保持不变。

可以证明，满足要求的修改结果存在且唯一。当 \(k=r-l\) 时，本次操作不改变任何元素。

\textbf{操作二：\texttt{2\ l\ r\ u\ v}，}对每个 \(l\le i\le r\)，令 \(a_i\leftarrow ua_i+v.\)

\textbf{操作三：\texttt{3\ l\ r}，}输出 \(\sum\limits_{i=l}^r a_i\) 对 \(P\) 取模后的值。

\problemheading{输入格式}{probB_input}

第一行包含一个整数 \(n\)。

第二行包含 \(n\) 个整数 \(a_1,a_2,\ldots,a_n\)。

第三行包含一个整数 \(m\)。

接下来 \(m\) 行，每行描述一次操作，格式如上。

\problemheading{输出格式}{probB_output}

对于每次操作三，输出一行一个整数，表示答案。答案应位于 \([0,P-1]\) 内。

\par\medskip\noindent\begin{minipage}{\linewidth}
\problemheading{样例 1 输入}{probB_sample_1_input}

\begin{verbatim}
8
1 4 6 9 13 18 24 31
8
3 1 8
1 1 8 3
3 1 8
3 4 7
2 3 6 2 1
3 1 8
1 2 7 2
3 5 8
\end{verbatim}
\end{minipage}\par\medskip

\par\medskip\noindent\begin{minipage}{\linewidth}
\problemheading{样例 1 输出}{probB_sample_1_output}

\begin{verbatim}
106
83
52
128
125
\end{verbatim}
\end{minipage}\par\medskip

\problemheading{样例 1 解释}{probB_sample_1_explanation}

第一次校准保留前四项，修改后的序列为

\[1,4,6,9,12,14,17,20.\]

长度为 \(3\) 的窗口和依次为 \(11,19,27,35,43,51\)，它们构成公差为 \(8\) 的等差数列。

仿射变换后的序列为 \(1,4,13,19,25,29,17,20\)。最后一次校准保留第 \(2\) 至第 \(4\) 项，校准后的整个序列为 \(1,4,13,19,28,34,43,20\)，所以最后一问的答案是 \(28+34+43+20=125\)。

\par\medskip\noindent\begin{minipage}{\linewidth}
\problemheading{样例 2 输入}{probB_sample_2_input}

\begin{verbatim}
5
4 7 2 9 1
10
3 1 1
2 3 3 0 10
3 2 4
1 1 3 2
3 1 5
1 1 5 1
3 3 5
2 1 1 2 3
2 5 5 0 0
3 1 5
\end{verbatim}
\end{minipage}\par\medskip

\par\medskip\noindent\begin{minipage}{\linewidth}
\problemheading{样例 2 输出}{probB_sample_2_output}

\begin{verbatim}
4
26
31
39
41
\end{verbatim}
\end{minipage}\par\medskip

\problemheading{样例 2 解释}{probB_sample_2_explanation}

第一次仿射变换把第 \(3\) 项设为 \(10\)。操作 \texttt{1\ 1\ 3\ 2} 的 \(k=r-l\)，因此序列仍为 \(4,7,10,9,1\)。

操作 \texttt{1\ 1\ 5\ 1} 后，长度为 \(1\) 的窗口和就是每一项，序列变为 \(4,7,10,13,16\)。最后两次单点仿射变换后，序列为 \(11,7,10,13,0\)。

\par\medskip\noindent\begin{minipage}{\linewidth}
\problemheading{样例 3 输入}{probB_sample_3_input}

\begin{verbatim}
6
998244351 1 998244351 4 0 3
12
3 1 6
1 1 5 2
3 1 6
2 2 5 998244352 998244352
3 1 6
2 1 6 0 0
1 2 6 1
3 2 5
2 6 6 998244352 998244352
3 5 6
1 4 6 1
3 1 6
\end{verbatim}
\end{minipage}\par\medskip

\par\medskip\noindent\begin{minipage}{\linewidth}
\problemheading{样例 3 输出}{probB_sample_3_output}

\begin{verbatim}
4
998244352
998244352
0
998244352
0
\end{verbatim}
\end{minipage}\par\medskip

\problemheading{样例 3 解释}{probB_sample_3_explanation}

在模 \(P\) 意义下，\(998244351=-2\)，\(998244352=-1\)。第一次校准后，序列为 \(-2,1,-2,1,-2,3\)，第 \(6\) 项保持为 \(3\)。随后的仿射变换 \(a_i\leftarrow-a_i-1\) 使序列变为 \(-2,-2,1,-2,1,3\)。

将整个序列设为 \(0\) 后，下一次校准仍得到全零序列。单点仿射变换把第 \(6\) 项改为 \(-1\)，最后一次校准又将它改回 \(0\)。

\problemheading{数据规模与约定}{probB_constraints}

对于所有数据，保证：

\begin{itemize}
\tightlist
\item
  \(1\le n,m\le10^5\)；
\item
  \(0\le a_i,u,v<P\)；
\item
  每次操作均满足 \(1\le l\le r\le n\)；
\item
  操作一还满足 \(1\le k\le r-l\)；
\item
  至少包含一次操作三。
\end{itemize}

本题采用分组计分。一个组内的所有测试点均通过，才能获得该组的分数。\\
下表中的 \(N=\max(n,m)\)。

\begin{centeredtable}\begin{tabularx}{0.85\textwidth}{|c|c|X|}\hline
组别 & 分值 & 额外限制\\\hline
1 & 20 & \(N\le1000\)\\\hline
2 & 15 & 不包含操作一\\\hline
3 & 15 & 包含操作一，且所有操作一均满足 \(k=1\)\\\hline
4 & 20 & \(N\le10000\)\\\hline
5 & 30 & 无\\\hline\end{tabularx}\end{centeredtable}
\newpage
\problemtitle{上下左右(lrud)}{probC_title}

\noindent 时间限制：1 秒\qquad 空间限制：512 MiB\par

\problemheading{题目描述}{probC_description}

给定正整数 \(N\)。

有一个 \(N \times N\) 的网格。从上往下数第 \(r\) 行、从左往右数第 \(c\) 列的格子记作 \((r, c)\)。

初始时，棋子放置在格子 \((1, 1)\) 上。

你需要恰好执行 \(N^2 - 1\) 次以下操作，使得棋子最终停在格子 \((N, N)\) 上：

\begin{itemize}
\tightlist
\item
  将棋子移动到当前所在格子的上下左右\textbf{或斜向}相邻格子之一。
\end{itemize}

并且，在整个移动过程中，\(N^2\) 个格子必须恰好各访问一次。初始放置棋子的格子 \((1, 1)\) 视为已被访问。

请计算一个操作序列中棋子恰好向右（不包括右上，右下）移动次数的最大值，并给出一组操作序列。

你需要解决 \(T\) 组这样的问题。

\problemheading{输入格式}{probC_input}

第一行一个整数 \(T\)，表示数据组数。

之后 \(T\) 行，每行一个整数 \(N\)。

\problemheading{输出格式}{probC_output}

对于每组数据，一行输出一个整数，表示向右移动次数的最大值；然后一行一个字符串表示操作序列，第 \(i\) 个字符表示路径中第 \(i\) 步的移动方向。其中，\texttt{Q}，\texttt{W}，\texttt{E}，\texttt{D}，\texttt{C}，\texttt{X}，\texttt{Z}，\texttt{A}分别代表左上，上，右上，右，右下，下，左下，左。

\par\medskip\noindent\begin{minipage}{\linewidth}
\problemheading{样例 1 输入}{probC_sample_1_input}

\begin{verbatim}
1
2
\end{verbatim}
\end{minipage}\par\medskip

\par\medskip\noindent\begin{minipage}{\linewidth}
\problemheading{样例 1 输出}{probC_sample_1_output}

\begin{verbatim}
2
DZD
\end{verbatim}
\end{minipage}\par\medskip

\problemheading{数据规模与约定}{probC_constraints}

所有数据均满足 \(1\le T\le5\)、\(2\le N\le1000\)。

本题采用分组计分。一个子任务内全部测试点通过后，才能获得该组分数；各组分数独立相加。

\begin{centeredtable}\begin{tabularx}{0.85\textwidth}{|c|c|X|}\hline
子任务 & 分值 & 额外限制\\\hline
1 & 10 & \(N\le3\)\\\hline
2 & 15 & \(N\le4\)\\\hline
3 & 20 & 每组数据满足 \(N\le4\) 或 \(N\) 为奇数\\\hline
4 & 20 & 每组数据满足 \(N\le4\) 或 \(N\not\equiv0\pmod4\)\\\hline
5 & 35 & 无额外限制\\\hline\end{tabularx}\end{centeredtable}

本题使用特殊评测。只要输出的路径合法且向右移动次数达到最大值，即可通过；操作序列不必与样例或标准答案相同。

\problemheading{提示}{probC_section_7}

本题附加文件下发了 \texttt{checker.cpp}，你可以用它来判断自己构造是否正确。不保证实际测评中使用同样的 checker.

使用方法：编译后在终端运行：\texttt{./checker\ {[}infile{]}\ {[}outfile{]}}（实际输入文件名时不带中括号）
\newpage
\problemtitle{存储空间(storage)}{probD_title}

\noindent 时间限制：10 秒\qquad 空间限制：256 MiB\par

\problemheading{题目描述}{probD_description}

给定一棵有根有序二叉树，共有 \(n\) 个叶子和 \(N=2n-1\) 个节点，根为节点 \(1\)。每个节点要么没有孩子，要么恰好有一个左孩子和一个右孩子。

将所有叶子按从左到右的顺序编号为 \(1,2,\ldots,n\)。这里的顺序是：递归遍历时，总是先遍历左子树，再遍历右子树。\textbf{叶子的顺序编号与输入中的节点编号是两个不同的概念。}

每个节点 \(v\) 对应一个区间 \(I_v=[L_v,R_v]\)，表示它的子树中所有叶子的顺序编号。节点 \(v\) 还有一个正整数存储费用 \(w_v\)。

记 \([l,r]=\{x\in\mathbb Z:l\le x\le r\}\)。你需要维护一个初始为空的区间集合 \(T\)，其中每个区间都是 \([1,n]\) 的非空子区间。

一个\textbf{存储方案}是树节点的一个集合 \(P\)。若对于每个 \(J\in T\)，都存在 \(U\subseteq P\)，使得 \(\{I_v:v\in U\}\) 中的区间两两无交，且它们的并集恰好为 \(J\)，则称 \(P\) 是当前 \(T\) 的一个可行存储方案。

方案 \(P\) 的总费用为 \(\sum\limits_{v\in P}w_v.\)

你需要处理 \(q\) 次操作：

\begin{itemize}
\tightlist
\item
  \texttt{+\ l\ r}：将区间 \([l,r]\) 加入 \(T\)；
\item
  \texttt{-\ l\ r}：将区间 \([l,r]\) 从 \(T\) 中删除。
\end{itemize}

保证每次加入的区间当前不在 \(T\) 中，每次删除的区间当前在 \(T\) 中。

每次操作后，输出当前可行存储方案的\textbf{最小总费用}，以及\textbf{达到该最小费用的不同存储方案数量}。方案数对 \(998244353\) 取模。

两个存储方案不同，当且仅当它们选取的节点集合不同。同一方案对某个目标区间的不同拼接方式，不额外计为不同方案。

注意，整个 \(P\) 允许同时包含祖先与子孙节点。两两无交的要求只针对覆盖某个目标时选出的子集 \(U\)；覆盖不同目标时，可以重复使用同一个存储节点。

\problemheading{输入格式}{probD_input}

第一行包含两个整数 \(n,q\)。

接下来 \(N=2n-1\) 行，第 \(v\) 行包含三个整数 \(\mathrm{ls}_v,\mathrm{rs}_v,w_v\)，分别表示节点 \(v\) 的左孩子编号、右孩子编号和存储费用。

若节点 \(v\) 是叶子，则 \(\mathrm{ls}_v=\mathrm{rs}_v=0\)；否则两者都是 \([1,N]\) 中的非零编号。除根为 \(1\) 外，节点编号没有其他顺序保证。

接下来 \(q\) 行，每行包含一个字符 \(\mathrm{op}\) 和两个整数 \(l,r\)。\(\mathrm{op}\) 为 \texttt{+} 或 \texttt{-}，表示一次操作。端点 \(l,r\) 使用叶子的顺序编号。

\problemheading{输出格式}{probD_output}

输出 \(q\) 行，每行包含两个非负整数，分别为本次操作后的最小总费用，以及最优存储方案数量对 \(998244353\) 取模的结果。

最小总费用不取模，可能超过 \(32\) 位有符号整数的范围。

\par\medskip\noindent\begin{minipage}{\linewidth}
\problemheading{样例 1 输入}{probD_sample_1_input}

\begin{verbatim}
2 6
2 3 2
0 0 1
0 0 1
+ 1 2
+ 1 1
+ 2 2
- 1 1
- 1 2
- 2 2
\end{verbatim}
\end{minipage}\par\medskip

\par\medskip\noindent\begin{minipage}{\linewidth}
\problemheading{样例 1 输出}{probD_sample_1_output}

\begin{verbatim}
2 2
2 1
2 1
2 1
1 1
0 1
\end{verbatim}
\end{minipage}\par\medskip

\problemheading{样例 1 解释}{probD_sample_1_explanation}

节点 \(1\) 对应 \([1,2]\)，费用为 \(2\)；节点 \(2,3\) 分别对应 \([1,1],[2,2]\)，费用均为 \(1\)。

第一次操作后，仅需覆盖 \([1,2]\)。选择根节点，或选择两个叶子，费用都为 \(2\)，所以共有两种最优方案。

第二次操作后还需覆盖 \([1,1]\)。选择两个叶子的总费用为 \(2\)；选择根和左叶子的费用为 \(3\)。因此只有前一种方案最优。

第三、四次操作后，选择两个叶子仍是唯一的最优方案。第五次操作后，仅剩右侧单点需求，最优费用为 \(1\)，方案数为 \(1\)。

最后一次操作后，\(T=\varnothing\)。空存储方案是唯一的最优方案，因此输出 \texttt{0\ 1}。

\problemheading{数据规模与约定}{probD_constraints}

对于所有数据，保证：

\begin{itemize}
\tightlist
\item
  \(1\le n,q\le 10^5\)；
\item
  \(N=2n-1\)，\(0\le\mathrm{ls}_v,\mathrm{rs}_v\le N\)，\(1\le w_v\le 10^9\)；
\item
  输入构成一棵以 \(1\) 为根的合法有序二叉树，全部节点均可从根到达，每个非叶节点恰有两个孩子，且恰有 \(n\) 个叶子；
\item
  每次操作均满足 \(1\le l\le r\le n\)，且所有操作合法。
\end{itemize}

本题采用独立子任务计分。只有通过某个子任务中的全部测试点，才能获得该子任务的分数。表中未列出的约束均遵循上述全局限制。

\begin{centeredtable}\begin{tabularx}{0.85\textwidth}{|c|c|X|}\hline
子任务 & 分值 & 额外约束\\\hline
1 & 5 & \(n\le7\)，\(q\le100\)\\\hline
2 & 10 & \(nq\le2\times10^6\)\\\hline
3 & 10 & 每个非叶节点均按区间中点分割，定义见下文\\\hline
4 & 5 & 所有叶子的深度均不超过 \(30\)，根的深度为 \(0\)\\\hline
5 & 10 & 每个操作区间都恰好等于某个节点的区间 \(I_v\)\\\hline
6 & 10 & 每个非叶节点的左孩子都是叶子\\\hline
7 & 10 & 每个非叶节点的右子树叶子数不少于左子树叶子数\\\hline
8 & 10 & 所有操作均为加入操作\\\hline
9 & 5 & 每个操作区间的长度不超过 \(2\)，即 \(r-l+1\le2\)\\\hline
10 & 25 & 无额外约束\\\hline\end{tabularx}\end{centeredtable}

子任务 3 中，``按区间中点分割''是指：对每个非叶节点 \(v\)，令 \(m=\lfloor(L_v+R_v)/2\rfloor\)，其左孩子的区间恰为 \([L_v,m]\)，右孩子的区间恰为 \([m+1,R_v]\)。
% END GENERATED PROBLEMS
\end{document}
